\appendix
\section{Proof and verification details}\label{app:proofs}

\subsection{Smooth active-set candidates}

On a branch \(F_z(x)=ax+b\), differentiating \eqref{eq:segmentpayoffs} gives
\[
4-b-a(p-q)-ap=0,
\qquad
a(p-q)+b-aq=0.
\]
Solving yields \eqref{eq:smoothcandidate}.  Applying it to the five positive-slope branches produces the following candidate differences:
\[
\begin{array}{c|c}
\text{branch} & x=p-q\\ \hline
-z\le x\le0 & 2(2-z)/3\\
0\le x\le z & 2(2-z)/9\\
z\le x\le1-z & 1/3\\
1-z\le x\le1 & 2(1+z)/9\\
1\le x\le1+z & 2(z-1)/3.
\end{array}
\]
For \(0<z<1/3\), the first lies above zero, the second lies above \(z\), the fourth lies below \(1-z\), and the fifth is negative.  Only \(x=1/3\) lies in its generating branch.

At a kink, a mutual local optimum would require the relevant own-price derivative to change sign in the correct direction.  For the poacher,
\[
\partial_q u_P=F_z-qF_z',
\]
and the slope-increase kinks \(x=-z,0,z\) fail that condition.  For the incumbent,
\[
\partial_p u_I=4-F_z-pF_z',
\]
and the slope-decrease kinks \(x=1-z,1,1+z\) fail the corresponding condition.  Flat-demand boundaries and zero prices do not generate mutual best responses.  This gives the pure-candidate exhaustiveness used in Proposition~\ref{prop:pure}.

\subsection{Selected mixed continuation}

At \(p^*\) in \eqref{eq:mixedprices}, \(q_L\) lies on the four-active branch and \(q_H\) on the three-active branch.  Direct substitution gives
\[
q_L\,4(p^*-q_L)
=
q_H\,[3(p^*-q_H)+z],
\]
so the poacher is indifferent between the two atoms.  The probability \(\lambda(z)\) in \eqref{eq:lambda} is obtained by setting the incumbent's expected marginal payoff to zero at \(p^*\).  Branch-by-branch comparison shows that neither player can improve outside the displayed support.  The project repository contains a second implementation that evaluates primitive clipped switching masses and globally maximizes deviations numerically; it does not call the symbolic branch solver.

\subsection{Backward-induction identities}

Let \(A_S,B_S\) denote the investor and non-investor gross payoffs at a one-investor sharing history, and \(A_{NS},B_{NS}\) the analogous no-sharing payoffs.  The reconstructed continuation payoffs satisfy
\[
A_S-M=F-B_S=K_S,
\]
and
\[
A_{NS}-M=N-B_{NS}=K_{NS}.
\]
The investment gain is therefore independent of the rival's Stage-II action within each symmetric Stage-I history, which yields the complete equality correspondences in Proposition~\ref{prop:stage2}.

For an asymmetric Stage-I history, the sharer's gain from investment is again \(K_S-k\) while the non-sharer's is \(K_{NS}-k\).  The identity
\[
K_{NS}-K_S=\frac{5z^2}{18}
\]
explains why the no-sharing side retains a strictly larger acquisition incentive throughout the maintained domain.

\subsection{Verification architecture}

The repository's regression suite includes:
\begin{enumerate}
\item an exact rational reproduction of the \(23/225\) deviation;
\item a symbolic active-set certificate;
\item an independent primitive-demand best-response evaluator;
\item a diagnostic two-by-two support search for alternative mixed equilibria;
\item a separate Stage-II/Stage-I correspondence audit;
\item direct numerical integration of consumer utility and welfare;
\item the two pre-specified portability falsification diagnostics; and
\item a Lean 4/mathlib proof-critical algebraic certificate.
\end{enumerate}
The numerical diagnostics are used as falsification or reconstruction checks where stated; they are not relabeled as analytic proofs.
